% =========================================================================
% Capítulo 5: Mecánica Lagrangiana
% Tema 01: Mecánica Lagrangiana
% Problema 08: Cilindro rodante en la cavidad de un bloque móvil
% =========================================================================
\subsubsection{Cilindro rodante en la cavidad de un bloque móvil}

Un bloque de masa $M$ descansa sobre una superficie horizontal perfectamente lisa (sin
fricción). El bloque posee una cavidad interior semicircular de radio $R$, cuyo centro de
curvatura es $O$, como se muestra en la Figura \ref{fig:c5_e01_p08_bloque_cavidad}. Dentro
de la cavidad, un cuerpo circular uniforme (cilindro sólido o esfera) de masa $m$ y radio
$r$ rueda sin resbalar sobre la superficie interior. El sistema está sujeto a la gravedad
$\mathbf{g}$ dirigida hacia abajo.

\begin{enumerate}
    \item Determine los grados de libertad y elija coordenadas generalizadas
          independientes ($x$ para el bloque y $\theta$ para la posición angular de $m$).
    \item Aplique la condición de rodadura sin deslizamiento para relacionar la rotación
          propia del cuerpo con $\theta$.
    \item Construya la función Lagrangiana $\mathcal{L}=T-V$ del sistema completo.
    \item Determine las ecuaciones diferenciales de movimiento para las coordenadas
          elegidas.
\end{enumerate}

\begin{figure}[!ht]
    \centering
    \begin{tikzpicture}[>=Latex, scale=1.1]

        % --- PARAMETROS ---
        \def\R{3}        % Radio de la cavidad del bloque M
        \def\r{0.8}      % Radio de la masa pequeña m

        % --- SUELO ---
        \fill[pattern=north east lines] (-4.5, -0.2) rectangle (5.5, 0);
        \draw[thick] (-4.5, 0) -- (5.5, 0);

        % --- BLOQUE M ---
        \coordinate (O) at (0, \R);
        \shade[left color=gray!40, right color=gray!10, middle color=gray!30]
            (-4, 0) -- (4, 0) -- (4, \R) -- (\R, \R) arc (0:-180:\R) -- (-4, \R) -- cycle;
        \draw[thick] (-4, 0) -- (4, 0) -- (4, \R) -- (\R, \R) arc (0:-180:\R) -- (-4, \R) -- cycle;
        \node at (-3.5, 0.4) {$M$};

        % --- CENTRO O ---
        \fill (O) circle (1.5pt) node[above left] {$O$};
        \draw[thick, dotted] (O) -- (0, 0.5);
        \draw[<->] (O) -- ++(210:\R) node[midway, above left] {$R$};

        % --- MASA m ---
        \path (O) ++(-50:{\R-\r}) coordinate (P);
        \path (O) ++(-50:\R) coordinate (Q);
        \shade[top color=gray!10, bottom color=gray!40] (P) circle (\r);
        \draw[thick] (P) circle (\r);
        \fill (P) circle (1pt) node[left=2pt] {$P$};
        \node at ([xshift=-12pt, yshift=-12pt]P) {$m$};
        \node at (Q) [above left=2pt] {$Q$};

        % --- LINEA OP Y ANGULO theta ---
        \draw[thick, dotted] (O) -- (P);
        \draw[->] (O) ++(-90:1.3) arc (-90:-50:1.3) node[midway, below] {$\theta$};

        % --- ANGULO phi Y VECTOR r ---
        \draw[thick, dotted] (P) -- ++(130:1);
        \draw[->] (P) -- ++(20:\r) node[midway, below] {$r$};
        \draw[->] (P) ++(130:0.5) arc (130:20:0.5) node[midway, below] {$\phi$};

        % --- GRAVEDAD ---
        \draw[->, thick] (-4.5, \R-0.5) -- (-4.5, \R-1.5) node[midway, left] {$\mathbf{g}$};

        % --- TRASLACION x ---
        \draw[|->] (0.5, -0.6) -- (1.5, -0.6) node[below] {$x$};

        % --- ANOTACIONES ---
        \node[font=\small] at (-1.7, 1.4) {No desliza};
        \draw (-1.7, 1.1) -- (0.5, 0.6);
        \node[font=\small] at (4.7, \R-0.5) {Sin fricción};
        \draw (4.7, \R-0.7) -- (3.8, 0.1);

    \end{tikzpicture}
    \caption{Cuerpo circular de masa $m$ que rueda sin resbalar en la cavidad semicircular de un bloque $M$ libre de moverse.}
    \label{fig:c5_e01_p08_bloque_cavidad}
\end{figure}

% --- SOLUCIÓN ---------------------------------------------------------------
\begin{solucion}{Cilindro rodante en la cavidad de un bloque móvil}
\subsubsection*{Pregunta 1--2: Vínculo cinemático y grados de libertad}

El sistema posee inicialmente tres coordenadas: $x$ (posición del bloque $M$), $\theta$
(posición angular del centro de $m$ respecto a $O$) y $\phi$ (rotación propia de $m$). La
rodadura sin deslizamiento vincula $\theta$ y $\phi$. Cuando el centro de $m$ avanza un
ángulo $\theta$, su centro recorre un arco $(R-r)\theta$ respecto al bloque; igualándolo al
arco de rotación propia $r\phi$:
$$r\phi = (R-r)\theta \;\Longrightarrow\; \dot\phi = \frac{R-r}{r}\dot\theta.$$
La velocidad angular absoluta del cuerpo es
$$\omega = \dot\theta + \dot\phi = \frac{R}{r}\dot\theta.$$
El sistema se reduce así a \textbf{dos grados de libertad}: $x$ y $\theta$.

\subsubsection*{Pregunta 3: Lagrangiano}

Sea $x$ la posición horizontal del centro $O$ de la cavidad. El bloque se mueve solo
horizontalmente ($v_M=\dot x$). El centro de masa de $m$ es
$$\vec{r}_m = \big(x + (R-r)\sin\theta\big)\hat{\imath} - (R-r)\cos\theta\,\hat{\jmath},$$
con
$$v_m^2 = \dot{x}^2 + (R-r)^2\dot\theta^2 + 2\dot{x}(R-r)\dot\theta\cos\theta.$$

Con $I=\gamma m r^2$ ($\gamma=\tfrac12$ cilindro, $\gamma=\tfrac25$ esfera), la energía
cinética es
$$T = \tfrac12 M\dot{x}^2 + \tfrac12 m v_m^2 + \tfrac12 I\omega^2
= \tfrac12(M+m)\dot{x}^2
+ \tfrac12\Big[m(R-r)^2 + I\tfrac{R^2}{r^2}\Big]\dot\theta^2
+ m(R-r)\dot{x}\dot\theta\cos\theta.$$

La energía potencial (nivel de referencia en $O$) es $V=-mg(R-r)\cos\theta$, de modo que
$$\boxed{\;\mathcal{L} = \tfrac12(M+m)\dot{x}^2
+ \tfrac12\Big[m(R-r)^2 + I\tfrac{R^2}{r^2}\Big]\dot\theta^2
+ m(R-r)\dot{x}\dot\theta\cos\theta + mg(R-r)\cos\theta\;}$$

\subsubsection*{Pregunta 4: Ecuaciones de movimiento}

La coordenada $x$ es cíclica ($\partial\mathcal{L}/\partial x=0$), por lo que se conserva
el momento conjugado $p_x=(M+m)\dot{x}+m(R-r)\dot\theta\cos\theta$. Derivando y aplicando
Euler-Lagrange en $\theta$:
$$\boxed{\begin{cases}
(M+m)\ddot{x} + m(R-r)\big(\ddot\theta\cos\theta - \dot\theta^2\sin\theta\big) = 0 \\[4pt]
\Big[m(R-r)^2 + I\left(\dfrac{R}{r}\right)^2\Big]\ddot\theta + m(R-r)\ddot{x}\cos\theta
+ mg(R-r)\sin\theta = 0
\end{cases}}$$
Se sustituye $I=\tfrac12 mr^2$ (cilindro) o $I=\tfrac25 mr^2$ (esfera) según corresponda.
\end{solucion}
